% 同一六模式族 V 上并排检验三种打散（行 = x1，列 = x2，次序 正常、广告、诈骗）。
% Natarajan：要求两行两列的四角都实心。V 中任意两行只共享一列 → 失败，d_N(V)=1。
% 图打散：存在参考格 f，使 f 格、f 行其余、f 列其余、其余四格四个区域各有实心点。
%   f=(正常,广告)：见证 (正常,广告)(正常,正常)(广告,广告)(诈骗,诈骗) → 成功，d_G(V)=2。
% DS：每个实心点在本行、本列各有另一实心点 → 成功，d_DS(V)=2。
\begingroup
\definecolor{LtSixGreen}{HTML}{829C8B}
\definecolor{LtSixGreenFill}{HTML}{E8EEE9}
\begin{tikzpicture}[figure picture,x=1mm,y=1mm]
 \def\cs{8}
 \newcommand{\ltsixbase}[1]{%
  \foreach \r in {1,2,3}{\foreach \c in {1,2,3}{
    \draw[FigureGrid,line width=.5pt] ({#1+\cs*(\c-1)},{16-\cs*(\r-1)}) rectangle ++(\cs,\cs);}}
  \foreach \lab [count=\c] in {正,广,诈}{
    \node[anchor=south,font=\figurefont\scriptsize] at ({#1+\cs*(\c-.5)},24.3) {\lab};}
  \foreach \lab [count=\r] in {正,广,诈}{
    \node[anchor=east,font=\figurefont\scriptsize] at ({#1-.6},{20-\cs*(\r-1)}) {\lab};}}
 \newcommand{\ltsixdots}[1]{%
  \foreach \r/\c in {1/1,1/2,2/2,2/3,3/3,3/1}{
    \fill[FigureBlue] ({#1+\cs*(\c-.5)},{20-\cs*(\r-1)}) circle (1.35);}
  \foreach \r/\c in {1/3,2/1,3/2}{
    \draw[FigureMuted,line width=.55pt,fill=white] ({#1+\cs*(\c-.5)},{20-\cs*(\r-1)}) circle (1);}}
 % ---------- Natarajan
 \node[figure heading] at (16,33) {Natarajan打散};
 \ltsixbase{4}
 \ltsixdots{4}
 \draw[FigureRed,line width=.8pt,dashed] ({4+\cs*.5},20) rectangle ({4+\cs*1.5},{20-\cs});
 \draw[FigureRed,line width=.8pt] ({4+\cs*.5-1.3},{20-\cs-1.3})--++(2.6,2.6)
   ({4+\cs*.5-1.3},{20-\cs+1.3})--++(2.6,-2.6);
 \node[figure note,anchor=north,align=center] at (16,-1.5) {固定两行两列，\\四角都要可实现};
 \node[anchor=north] at (16,-10) {$d_N(\mathcal V)=1$};
 % ---------- graph
 \node[figure heading] at (55,33) {图打散};
 \fill[FigureYellowFill] ({43+\cs},16) rectangle ++(\cs,\cs);
 \fill[FigureBlueFill] (43,16) rectangle ++(\cs,\cs);
 \fill[FigureBlueFill] ({43+2*\cs},16) rectangle ++(\cs,\cs);
 \fill[LtSixGreenFill] ({43+\cs},{16-2*\cs}) rectangle ++(\cs,{2*\cs});
 \fill[FigureRedFill] (43,{16-2*\cs}) rectangle ++(\cs,{2*\cs});
 \fill[FigureRedFill] ({43+2*\cs},{16-2*\cs}) rectangle ++(\cs,{2*\cs});
 \ltsixbase{43}
 \ltsixdots{43}
 \foreach \r/\c in {1/2,1/1,2/2,3/3}{
   \draw[FigureInk,line width=.65pt] ({43+\cs*(\c-.5)},{20-\cs*(\r-1)}) circle (2.2);}
 \node[figure note,anchor=north,align=center] at (55,-1.5) {选一个参考格，\\四个区域各有一点};
 \node[anchor=north] at (55,-10) {$d_G(\mathcal V)=2$};
 % ---------- DS
 \node[figure heading] at (94,33) {DS打散};
 \ltsixbase{82}
 \foreach \r/\c/\n in {1/1/a,1/2/b,2/2/c,2/3/d,3/3/e,3/1/f}{
   \coordinate (q\n) at ({82+\cs*(\c-.5)},{20-\cs*(\r-1)});}
 \draw[FigureBlue,line width=.9pt] (qb)--(qc) (qd)--(qe);
 \draw[FigureBlue,line width=.9pt] (qa) to[bend right=35] (qf);
 \draw[LtSixGreen,line width=.9pt,dash pattern=on 2pt off 1.2pt] (qa)--(qb) (qc)--(qd);
 \draw[LtSixGreen,line width=.9pt,dash pattern=on 2pt off 1.2pt] (qf) to[bend right=30] (qe);
 \ltsixdots{82}
 \node[figure note,anchor=north,align=center] at (94,-1.5) {每个点在本行、\\本列都有邻居};
 \node[anchor=north] at (94,-10) {$d_{DS}(\mathcal V)=2$};
\end{tikzpicture}
\endgroup
